\chapter{数字校准的完整数学模型与硬件代价}
\label{ch:cal-depth}
本章回答三个不同问题：物理开关状态怎样成为重构公式，公式怎样成为有精确异常语义的整数运算，以及这些运算为什么可能成为综合瓶颈。所用 Q30、Q32、96-bit 受检域及 63-bit 除法分子均来自本工程的冻结算术契约，不能反向归因为论文公布的芯片内部位宽。以下推导主要对应 \code{cal_weight_reduce}、\code{cal_residue_mac}、\code{adc2_dec}、\code{div_floor}、\code{cal_output_stage} 和 \code{recon_core}。

\section{从物理单元开始，而不是从理想码开始}
设物理单元由二元索引 \((s,u)\) 标识，\(s\) 为 slice，\(u\) 为该 slice 内的电容单位。本工程生产几何有 18 个 slice，每个含 63 个主组和 8 个子组单位，共 1278 个系数。一次转换的活动集合记为 \(\mathcal A_n\)，正常情况下含 8 个不同的物理 slice。权重 \(w_{s,u}\) 是配置期载入的正数，包含所选择的等效归一化和物理非理想校正信息。DEM 改变参与本次转换的物理地址和开关状态，\emph{不会}把系数的物理身份一起重新编号。

令 \(b_{n,s,u}\in\{0,1\}\) 为本次实际锁存的底板选择，\(a_{n,s,u}\in\{0,1\}\) 表示该单位是否参与输入采样增益，\(\rho_{n,s,u}\) 表示采样 dither 的额外轨贡献。用下面三个量保留开关网络对重构的全部数字信息：
\begin{align}
T_n &= \sum_{s\in\mathcal A_n}\sum_u w_{s,u},\\
G_n &= \sum_{s\in\mathcal A_n}\sum_u a_{n,s,u}w_{s,u},\\
R_n &= \sum_{s\in\mathcal A_n}\sum_u
       (1-2b_{n,s,u}+\rho_{n,s,u})w_{s,u}.
\end{align}
普通采样时 \(a=1,\rho=0\)，因此 \(G=T\)。在当前采样掩码模型中，少数保留单位 \(a=0\)，它们从采样增益中扣除，并以 \(\rho=\pm1\) 进入轨贡献。这里描述的是工程模型的开关系数；最终电荷方程仍须由实际电容连接、参考极性和采样时序核验。若模拟宏采用不同的底板极性，必须同时修改接口契约和 oracle，不能只在最终输出上补一个符号。

在统一归一化电压域，可将重构关系写成
\begin{equation}
\hat x_n=\frac{F_n-O-R_n-T_n I_n}{G_n}.
\label{eq:physical-calibration}
\end{equation}
其中 \(F_n\) 是后级码对应的电压估计，\(O\) 是偏移校正，\(I_n\) 是已知公共注入电压。式中隐含权重、电压量纲已经归一到相同基准；权重可按指定电容基准定义为 \(w=C_{\rm eff}/C_{\rm f}\)，其基准及增益网络必须固定；若直接使用伏特和法拉，则须重新列写各节点的电荷方程，不能在此归一化式中直接代入法拉数值。RTL 消费的是冻结的整数编码，不能把任意单位的模拟数据直接填入这些端口。

\subsection{为什么只按逻辑码查一张表不够}
考虑两个名义相同、实际权重分别为 \(1+\epsilon\) 与 \(1-\epsilon\) 的物理单位。两个样本都要求“接通一个单位”，理想码相同，但 DEM 分别选中这两个单位时，\(R\) 相差 \(4\epsilon\)。若校正仅依赖逻辑码而丢弃物理选择，这个误差会被当成随机噪声保留；物理加权则能在模型成立时逐次扣除它。

因此，必须先完成物理地址译码，再根据同一份开关快照求和。系数表按物理地址排列，开关掩码按本次活动槽位排列，这两种索引空间不可混用。系统中的 sample ID 只能证明事务标签一致，不能单独证明物理地址排列正确；需要逐项系数装载、已知单元扰动和地址置换负控共同构成证据。

\begin{figure}[htbp]\centering
\begin{tikzpicture}[node distance=8mm, every node/.style={font=\small}, box/.style={draw,rounded corners,align=center,minimum height=12mm,text width=31mm},>=Stealth]
\node[box] (phys) {锁存的物理地址\\与开关状态};
\node[box,right=of phys] (coef) {物理系数对齐\\1278 项配置};
\node[box,right=of coef] (terms) {加权归约\\\(T,G,R\)};
\node[box,right=of terms] (num) {电压重构\\精确整数运算};
\draw[->] (phys)--(coef);\draw[->] (coef)--(terms);\draw[->] (terms)--(num);
\node[box,below=of terms] (fine) {同一残差样本\\\(F,O,I\) 与 ID};
\draw[->] (fine)-|(num);
\end{tikzpicture}
\caption{本工程校准数据流。图为依据源码绘制的解释图，不是芯片内部原理图或测量波形。}
\end{figure}

\section{定点格式、缩放顺序与单位闭合}
定义 \(Q_w=2^{30}\)、\(Q_v=2^{32}\)，物理权重整数 \(W=\operatorname{round}(Q_ww)\)，归一化电压整数 \(V=\operatorname{round}(Q_vv)\)。求和结果 \(T,G,R\) 在下文改用整数形式，仍具有权重 Q30 的比例因子。相应整数分子为
\begin{equation}
N=(F-O)Q_w-RQ_v-TI.
\end{equation}
三项均处于 Q62 的乘积尺度。若错误地用 \(R\) 直接减去 Q32 电压量，数值可能在理想零点附近凑巧接近，但随输入和 DEM 改变会产生系统性比例错误。逐项标注单位是发现这一类问题最直接的方法。

20-bit offset-binary 输出把 \([-1,1)\) 映射到 \([0,2^{20})\)。因此，在还没有饱和处理前，
\begin{align}
C_{\rm raw}
&=\left\lfloor\frac{\hat x+1}{2}\,2^{20}\right\rfloor\\
&=\left\lfloor\frac{(N+G2^{32})2^{20}}{G2^{33}}\right\rfloor.
\end{align}
最后输出 \(C=\min(2^{20}-1,\max(0,C_{\rm raw}))\)，同时报告高、低剪裁。饱和映射和计算失败是两个不同结果：饱和是合法数学结果超出码域；累加溢出或增益错误意味着本次数学运算不可信，应保留上一笔可用码并携带错误标志。

\begin{table}[htbp]\centering\small
\begin{tabularx}{\textwidth}{p{29mm}p{29mm}X}\toprule
信号/常量&整数尺度&工程含义与误用风险\\\midrule
\(W,T,G,R\)&Q30 权重&求和结果不能按 Q32 电压直接解释。\\
\(F,O,I\)&Q32 电压&相对于同一归一化满量程；不是原始 ADC2 无符号码。\\
\(N\)&Q62 乘积&减法前需显式符号/零扩展，不能依赖上下文隐式推宽。\\
\(G2^{33}\)&正整数分母&同时包括电压尺度和双极到单极的系数 2。\\
\(2^{20}\)&输出计数尺度&来自工程输出码宽，不能由 ADC2 码宽替换。\\
\bottomrule\end{tabularx}
\caption{重构公式的量纲检查表。}
\end{table}

\subsection{二次 floor 等价为何成立}
为减少除法输入宽度，RTL 先令 \(S=(N+G2^{32})2^{20}\)，再计算
\begin{equation}
A=\left\lfloor S/2^{33}\right\rfloor,
\qquad C_{\rm raw}=\lfloor A/G\rfloor.
\end{equation}
对任意整数 \(S\) 及正整数 \(m,G\)，写成欧几里得分解 \(S=mA+r\)，其中 \(0\le r<m\)，再令 \(A=Gq+t\)，其中 \(0\le t<G\)，便有
\begin{equation}
S=mGq+mt+r,\qquad 0\le mt+r<mG.
\end{equation}
所以 \(\lfloor S/(mG)\rfloor=q=\lfloor\lfloor S/m\rfloor/G\rfloor\)。该证明包含负 \(S\)，前提是两次都向负无穷取整。算术右移在正确有符号表达式上给出除以二次幂的 floor；SystemVerilog 的有符号除法则向零截断，不能直接代换最后一步。

例如 \(S=-17,m=4,G=3\)，正确结果为 \(\lfloor-17/12\rfloor=-2\)。先 floor 得 \(-5\)，再 floor 除以 3 仍为 \(-2\)；若错误地让第二步向零截断，则输出 \(-1\)。这一个码的偏差恰好出现在负侧边界，单独测试正弦中段或正斜坡不一定能触发。

\section{ADC2 码解码与半值向偶}
后级原始码 \(c\) 的范围是 \(0\) 到 \(2^{B_2}-1\)，其区间端点由配置 \(F_{\min},F_{\max}\) 给出。当前模型使用量化区间中点：
\begin{equation}
F=F_{\min}+\operatorname{round}_{\rm even}
\left(\frac{(2c+1)(F_{\max}-F_{\min})}{2^{B_2+1}}\right).
\end{equation}
中点模型的目的是不把码下边界当作无偏的残差估计；它本身不证明 ADC2 的实际传输曲线均匀。对非线性或非均匀 ADC2，需要更细的模型和标定数据，不能继续把一对端点当作完整校正。

令乘积 \(P=(2c+1)(F_{\max}-F_{\min})\)、分母 \(D=2^{B_2+1}\)，先求非负整数商 \(q=\lfloor P/D\rfloor\) 和余数 \(r=P-qD\)。舍入决策为
\begin{equation}
q'=q+\begin{cases}
1,&2r>D,\\
1,&2r=D\ \text{且}\ q\text{为奇数},\\
0,&\text{其他情况}.
\end{cases}
\end{equation}
只有等号分支区分 half-up 和 ties-to-even。独立 oracle 必须主动构造这些分支，随机浮点输入未必精确命中半值。范围差应在足够宽的有符号域内计算，不能先以 64-bit 减法回绕后再扩展。无效范围需要产生可追踪异常，而不是悄悄交换端点。

\subsection{解码误差怎样传播到输出}
固定其他项时，\(\delta\hat x=\delta F/G\)，所以残差码误差会被数字增益归一反映到最终输出。对小扰动，包含系数误差的一阶关系为
\begin{equation}
\delta\hat x\approx
\frac{\delta F-\delta O-\delta R-I\delta T-T\delta I}{G}
-\hat x\frac{\delta G}{G}.
\end{equation}
这说明为什么分母误差不可只视为一个固定偏移：其影响随输入幅度改变；如果活动集合变化，\(\delta G\) 还随 DEM 变化。该公式可用于仿真预算，但不包含大扰动、饱和、比较器误判和动态建立误差，不能作为真实 ADC 的完整误差模型。

\section{位宽证明与合法范围}
\subsection{系数表和归约位宽}
载入接口要求 \(0<W<2^{47}\)。默认物理系数数目为 1278，故全表和满足
\begin{equation}
\sum W<1278\,2^{47}<2^{58}.
\end{equation}
即使按存储器支持的最大几何 \(32\times128\) 推算，也有 \(\sum W<2^{59}<2^{60}\)。因此在当前参数约束下，总和容量门禁是静态可证明为真，写入时不必形成一条读取旧系数、减去旧值、加上新值再作 64-bit 比较的运行路径。源码保留条件化后备检查，综合器可在静态条件成立时裁去它。这类消除来自数学边界，和删掉保护逻辑碰碰运气不同。

正常活动集合只有 8 个 slice，\(T\) 的实际最大项数是 \(8\times71=568\)。归约模块还要对直接输入的任意 48-bit 位型保持确定行为，以支持边界 miter；其参数守卫按完整物理项数加 \(\lceil\log_2N\rceil\) 位，默认 64-bit 仍有余量。配置合法域与叶模块测试域应分别说明，不宜用合法配置的更紧边界去解释非法位型测试。

\subsection{中间乘加为什么比输出宽得多}
电压输入为 64-bit 有符号，两个电压相减需要 65-bit。权重总和为无符号，乘以有符号注入量之前必须零扩展并显式转成 signed，否则一个无符号操作数就可能使整个表达式按无符号规则求值。\code{cal_residue_mac} 采用 130-bit 乘积容器、132-bit 运算容器，并分别检查关键项落在
\begin{equation}
-2^{95}\le z<2^{95}
\end{equation}
这一受检域内。容器宽度保证能观察溢出，而不是让越界先回绕成一个看似合法的值。

受检上界采用严格小于 \(2^{95}\)，下界包含 \(-2^{95}\)。对二进制补码，这种不对称正是类型范围的真实结构。把“绝对值小于上界”直接写成对称比较，会错误拒绝合法最小负值；现有负控特意修改了这一判断，独立测试确实能抓到。

\begin{longtable}{p{29mm}p{29mm}p{88mm}}
\toprule 阶段&实现容器&为什么需要这一范围\\\midrule\endfirsthead
\toprule 阶段&实现容器&为什么需要这一范围\\\midrule\endhead
\(F-O\)&65-bit signed&两个 64-bit 有符号端点相减可能超过原输入位宽。\\
\(TI\)&130-bit signed&先为无符号总和留符号位，再与有符号注入相乘；完整积保留溢出证据。\\
各乘加项及 \(N\)&132-bit signed&使各项先精确计算，再检查 96-bit 合同，而非先截断后检测。\\
\(N+G2^{32}\)&97-bit signed&在先前受检项合法时，容纳两个合法大数的中间和。\\
左移 20-bit 后 \(S\)&117-bit signed&左移必须先扩展；在窄容器内移位会丢失高位证据。\\
右移后 \(A\)&63-bit signed&只有 \(S\) 范围检查通过时，提取的 63-bit 数才可用于合法结果。\\
\bottomrule\end{longtable}

\subsection{异常路径上的数值是否还必须精确}
如果某个受检项已经越界，后续窄路径可以只维持确定的控制流程，因为可见行为规定为“报告本次失败并保留旧码”。但它必须满足两个条件：错误位属于本次事务，且不能被后续算术异常覆盖或清掉。反过来，只要所有受检条件都合法，输出就必须等于独立整数契约，不允许以“最终只看 20 bit”为理由提前截断。

这是一条有实际 PPA 意义的边界：安全的优化可以缩小仅在错误路径上使用的无意义计算，也可以提前旁路已经确定的饱和值；但应通过可观察语义证明，而不是要求错误输入有任意输出。对用户而言，码值、valid、sample ID、当前错误位和粘滞错误位都是接口的一部分。

\section{列共享归约：代数变化与电路变化}
当前采样 dither 对各活动 slice 的同一列采用相同轨符号。对第 \(j\) 个保留单位定义
\begin{equation}
C_j=\sum_{s\in\mathcal A}W_{s,u_j},\quad
M=\sum_jC_j,\quad
R_{\rm dith}=e_{\rm sm}\sum_j\sigma_jC_j,\quad \sigma_j\in\{-1,+1\}.
\end{equation}
其中 \(e_{\rm sm}\) 为采样掩码使能，关闭时 dither 轨贡献严格为零。则
\begin{equation}
G=\begin{cases}T-M,&\text{采样掩码开启},\\T,&\text{否则},\end{cases}
\qquad R=T-2W_{\rm on}+R_{\rm dith}.
\end{equation}
旧结构把条件取负放在每个物理单元后，再把带符号叶子相加；新结构先对同列的正权重求和，再对整列取正负。这是利用共享轨符号提取公因子，既没有近似，也不需要在线训练或额外配置状态。

对默认 18 个 slice、4 个保留列，条件取负从 72 处变为 4 处。保守扣除旧版可以共同消去的部分后，相关加减表达式由 249 项降为 75 项，差为 174 项。该计数仅用于解释硬件结构意图；综合器是否共享、映射到何种 carry/DSP、控制扇出如何变化，必须由功能有效的网表比较。特别是 \(G=T-M\) 引入最终减法，可能使增益路径变长，因此逻辑表达式减少并不自动意味着最高频率提高。

\begin{figure}[htbp]\centering
\begin{tikzpicture}[font=\small,>=Stealth,b/.style={draw,align=center,text width=38mm,minimum height=12mm},node distance=7mm]
\node[b] (a) {旧：各单元取正负\\\(\sigma_jW_{s,j}\)};
\node[b,right=of a] (b) {全部带符号叶子\\进入归约树};
\node[b,right=of b] (c) {独立形成增益\\与 dither 和};
\draw[->] (a)--(b);\draw[->] (b)--(c);
\node[b,below=13mm of a] (d) {新：先按列求和\\\(C_j=\sum_sW_{s,j}\)};
\node[b,right=of d] (e) {每列一次取正负\\\(\sigma_j C_j\)};
\node[b,right=of e] (f) {共享列和\\\(G=T-\sum_jC_j\)};
\draw[->] (d)--(e);\draw[->] (e)--(f);
\end{tikzpicture}
\caption{列共享改写的结构比较。两条路径的输出数学等价，实际面积和时序必须通过综合后的语义检查再衡量。}
\end{figure}

\subsection{为什么不能在这里偷偷假定 ID 永不重复}
正常调度应保证 8 个物理 ID 不同，但归约叶模块仍能被其他测试或调用方直接驱动。当前结构把逻辑掩码移动到物理行，同一物理行被重复指定时采用活动标记与掩码并集，并同时拉高 \code{invalid_slice}。它不是把同一权重重复累加两次。新旧比较必须保留这种确定的非法输入语义，否则 miter 只在正常 ID 上通过仍可能改变错误路径的状态。

这也解释了 miter 枚举集合与随机有序排列各自的作用：集合枚举覆盖全部 8-of-18 物理参与关系，排列测试覆盖逻辑槽位向物理行移动的译码关系。两种检查互补。\(\binom{18}{8}=43758\) 不能直接说成 18 个通道的所有时间序列；跨样本状态仍由协议测试承担。

\section{配置期与转换期的硬件分工}
\subsection{当前配置协议的真实性能含义}
\code{weight_store} 保存 61,344 bit 原始系数数据，并维护每个物理地址本 epoch 是否写过的 bitmap。\code{calib_regs} 维护三个标量的完整性和控制位。reset 或 clear 开启一个新的装载 epoch；只有所有必填项齐备、范围合法、没有同时写入及在途冲突，validate 才能提交配置。转换期拒绝权重和标量改写，因此一个样本不会混用两个系数版本。

61,344 bit 是逻辑状态容量，不是综合后触发器数量的精确承诺。合法正权重的最高位可能被优化，bitmap 和其他寄存器又会增加状态数；数组是否映射成 BRAM 还受读端口形式与复位语义约束。当前全表同拍可见接口天然要求大量并行读，单靠把声明换成 \code{logic mem [...]} 不会得到既省资源又行为相同的双口 RAM。

\subsection{可预计算项与不可预计算项}
配置期可以预计算每个 slice 的全行和、保留列和等与实时开关无关的量。转换时 \(T\) 可从 8 个行和归约，\(M\) 可从各活动 slice 的保留部分和归约，从而避免反复展开固定的局部求和。然而 \(W_{\rm on}\) 取决于本次 DEM 后的实际开关，每个任意系数组合仍需处理，不能仅以行和替代。

更激进的前缀和缓存只有在开关选择可表示为少量有序区间时才有收益。二维轮换、子组交换与不同物理排列会改变区间结构；如果为每种 DEM 状态都缓存全表校正，缓存容量、更新时间和地址译码又可能超过原来的加法树。应先从真实掩码轨迹统计区间数和复用度，再决定硬件形态。

\begin{table}[htbp]\centering\small
\begin{tabularx}{\textwidth}{p{31mm}X X}\toprule
候选结构&可能收益&必须支付的代价或证明\\\midrule
配置期行和/列和&减少每次重复相加的固定项&替换写入和完整重载时保持缓存与系数一致。\\
多 bank 系数存储&限制长距离布线、支持分拍读取&银行冲突、读口数量和 16 拍预算需要实际调度。\\
前缀和缓存&连续区间掩码变成少数相减&实际 DEM 置换必须仍能分解成有限区间。\\
固定镜像 ROM&大幅减少可写状态&失去片外拟合后更新能力，不符合当前通用校准接口目标。\\
压缩系数位宽&减少存储和加法器宽度&建立量化误差预算并重新验证全部系数与边界。\\
\bottomrule\end{tabularx}
\caption{后续优化的可审查取舍；本表是方案比较，不代表这些结构已经实现。}
\end{table}

\section{离线辨识与在线应用必须分别验收}
论文的片外系数获取与片上校正划分，允许 RTL 只实现确定的算术应用，并不要求在 RTL 中增加 LMS/RLS。未知权重的获取可写成线性或近似线性观测问题：\(\mathbf y=\mathbf H\boldsymbol\theta+\boldsymbol\eta\)。\(\boldsymbol\theta\) 包含待求的物理权重及必要偏移，\(\mathbf H\) 来自可观察的开关/激励状态，\(\boldsymbol\eta\) 包含量化和模型误差。

若 \(\mathbf H\) 秩不足，则不同系数组合能给出相同训练输出。即使数值优化器报告收敛，也不等于找到了真实物理权重。若列近似相关，条件数很大，极小的观测噪声就会使拟合系数变化；这种情况下应限制可辨识的参数组合、增加具有区分力的激励或加入有物理依据的约束，而不能仅增加迭代次数。

物理权重和电压归一化之间可能存在比例自由度。必须固定参考尺度或选定一个规范化约束，否则 \(G,R,T\) 一起缩放与后级幅度缩放会产生参数歧义。报告应说明固定了什么量、为什么该约束不吸收目标误差，以及每个系数的单位。仅发布最终系数矩阵无法复现辨识过程。

\subsection{留出验证究竟独立在哪里}
本工程的 128 笔拟合留出验证证明：已生成的 1278 项镜像确实写入当前 RTL，随后观测到的码、flags、ID 和延迟与给定期望一致。这是接口和整数应用的强证据。它不自动证明训练激励覆盖真实未知芯片的全部物理误差，也不证明留出数据与训练数据在温度、动态历史和噪声实现上独立。

进一步实验应分开四组数据：拟合用样本、同条件未参与拟合的留出样本、不同输入频率/幅度的泛化样本，以及不同温度/电源/历史条件的漂移样本。每组都保留原始残差码和开关状态，使失败时能重新执行校正而不重跑模拟。随机种子必须登记；同一个种子生成输入和模拟误差可能让模型误差被训练过程意外学习。

\subsection{对量化误差建立可计算预算}
令 \(s_i=1-2b_i+\rho_i\) 为完整轨系数。采用最近舍入且 \(|\delta w_i|\le 2^{-31}\)，最坏绝对和给出 \(|\delta T|\le N2^{-31}\)、\(|\delta R|\le\sum_i|s_i|2^{-31}\)。代入一阶误差式可筛查 Q30 是否足够；\(G\) 接近零会放大误差，合法系数和采样结构须保证其正下界。

该上界保守，因为不同权重量化误差可能抵消，但不能据此随意假设独立白噪声。用于缩位的正式依据应同时包含最坏误差界、真实系数分布下的位精确仿真以及最差输入/开关组合。任何缩位导致的输出码变化都需要在规范中明确，而不是当作“综合优化等价”提交。

\section{样本快照与输出协议的因果关系}
前端当前正在形成的 RDAC 残差与后级 ADC2 正在量化的残差可能属于不同样本。\code{cal_sample_context} 为二者保存独立上下文；residue capture 锁存当前物理选择、注入和 ID，quiet sample 再把上一份残差上下文转给后级。这一层的核心要求不是标签数字漂亮，而是所有校正量具有相同采样时刻和物理来源。

\code{recon_core} 在接受 start 时快照 \(T,G,R,F,O,I\) 和 sample ID，下一拍才能对这些寄存器计算溢出及增益错误。若在接受 start 的同一个 sequential 分支中读取这些寄存器来设置错误位，非阻塞赋值语义会使用上一样本的数据，导致异常晚一笔。大规模无异常斜坡不会区分这两个实现；只有正常、异常、正常交替且 ID 逐笔核对的测试能直接暴露。

输出 \code{dout_valid} 表示“本次事务完成”，并非无条件声明 \code{dout} 是本次的新有效测量。累加溢出或增益错误使输出保持上一笔码，同时给出本次 ID 和 result flags；单独出现 ADC2 范围告警时，输出仍按候选码更新并携带告警位。消费端必须区分这两种情况，不能把所有错误位都解释成旧码保持。若应用更需要 \code{sample_good}，可以从当前 flags 派生，不能把粘滞错误位直接当作当前样本的质量标签。

\begin{longtable}{>{\raggedright\arraybackslash}p{27mm}>{\raggedright\arraybackslash}p{40mm}>{\raggedright\arraybackslash}p{78mm}}
\toprule 情形&可见输出&消费端应做的处理\\\midrule\endfirsthead
\toprule 情形&可见输出&消费端应做的处理\\\midrule\endhead
正常计算&新码、当前 ID、valid、正常 flags&接受并按当前 ID 对齐。\\
合法剪裁&边界码、clip 标志、当前 ID&按量程溢出处理；不要解释为算术回绕。\\
累加/增益错误&旧码保持、当前事务 valid 和错误位&登记本次测量无效，不能把旧码赋予新样本。\\
配置撤销&输出与 valid 按 cfg 规则清除&结束当前配置 epoch，重新装载后再采纳样本。\\
ADC2 范围告警&按已定义输出规则更新并携带 ADC2 位&由应用质量策略决定是否采用，不能丢弃该位。\\
\bottomrule\end{longtable}

\section{校准功能成功的最小证据链}
完整的数字校准证据至少要连起五个环节：物理地址和系数镜像对应；实际开关快照对应当前残差；整数公式与独立 oracle 一致；异常及 ID 不串样本；综合映射后这些函数仍然存在。当前仿真提供了前四类的多个有界证据，而本轮 Vivado 的无驱动发现说明第五类不可省略。

每幅图应登记脚本、原始日志、源码SHA及采样条件。曲线显示趋势，逐笔比较判定成功，负控说明区分力。同一数据重画多图不会增加独立证据，应明确错误实例与RTL修复的因果关系。
